\section{Why Episodic Agents Stall}
\label{sec:episodic}

\begin{designbox}
Diagnose why a capable model, run long, degrades into busy sparsity --- and derive
the design requirements that a system must meet to make research \emph{compound}
across runs instead of restarting from zero.
\end{designbox}

\subsection{Model capability is not system capability}

It is easy to assume that a more capable model is, by itself, a more capable
research system. A stronger model does raise the ceiling of any single step: it
reads code more accurately, writes better prose, and reasons further in one call.
But a research \emph{system} is judged over thousands of steps sustained across
hours to days, and that horizon is governed by properties the weights do not
supply: whether the thread of work survives the end of a call, whether claims are
checked before they are believed, whether what was learned is retained and reused.
Model capability is per-call; system capability is cross-call. Treating them as the
same thing is why ``give a big model a hard problem and wait'' fails --- the missing
capability lives in the substrate around the model, not in the model.

\subsection{Episodic work pays context-recovery cost}

Left to run long, a model works in episodes. Each episode accumulates context
until the window overflows and the transcript is compacted or dropped; the next
episode then reconstructs where it was --- re-reading the same files, re-deriving
decisions it already made, re-discovering dead ends it already ruled out. This
context-recovery cost is pure overhead: activity with near-zero decision
efficiency. It is not a prompt-quality problem that a better instruction removes;
it is a structural property of unbounded session growth. The longer the horizon,
the larger the fraction of elapsed time the system spends paying this tax instead
of advancing the frontier.

\subsection{Unbounded context is not institutional memory}

The obvious response --- never drop context, just make the window bigger --- solves
the wrong problem. An ever-growing transcript is not memory. Institutional memory
is curated, addressable, and reusable: a specific technique that worked, a baseline
that was rejected and why, a verifier that caught a fake result. A monolithic
transcript is none of these; it is an undifferentiated log in which the few
reusable facts are buried among the reasoning that produced them, and it is
precisely what lossy compaction shreds first when the window finally overflows.
Durable capability has to live \emph{outside} the model, in persisted state that
can be written deliberately, retrieved selectively, and carried into a future run.

\subsection{Process data disappears when only the final artifact survives}

Most of the value of a research run is in its process, not its final artifact. The
manuscript or merged patch records the destination; it discards the map of how the
work got there --- what was attempted and failed, which measurement was found
contaminated, why one design was chosen over another. When a system keeps only the
final artifact, that process knowledge is thrown away at the end of every run, so
the next run cannot stand on it and is free to repeat the same mistakes. A system
that is meant to compound must persist the trajectory and the verdicts, not just
the deliverable, and must treat those process records as first-class, auditable
state.

\subsection{Design requirements for compounding research}

These failure modes are not independent bugs; together they define what a
long-horizon research system must provide. Five requirements follow directly, and
each is met by a concrete mechanism in the architecture that Act~II documents
(Section~\ref{sec:architecture}).

\begin{enumerate}[leftmargin=1.5em, label=\textbf{R\arabic*.}]
  \item \textbf{Persistent objective.} A durable, operator-anchored objective that
    survives restarts and upgrades, so the system always knows what it is pursuing
    and never silently re-plans a campaign from a stale premise.
  \item \textbf{Role separation.} Distinct model-driven roles for decision,
    execution, and verification, so no single stateless call owns all three and
    completion is never a self-report from the role that did the work.
  \item \textbf{Bounded context.} Per-session context held to an explicit bound and
    re-seeded from curated working memory, so the runaway compaction that drives
    context-recovery cost cannot occur.
  \item \textbf{Independent evidence.} Progress and completion judged by an
    authority that inspects artifacts and execution logs against evidence, rather
    than inferred from activity, filenames, or token volume.
  \item \textbf{Reusable runtime state.} Verified evidence deposited as durable,
    addressable capability --- memory, skills, tools, verifiers, routing,
    evaluations --- that the next unknown task inherits instead of starting from
    zero.
\end{enumerate}

Act~II turns these five requirements into a running system: a three-plane
architecture that separates decision authority, execution, and evidence, and the
runtime that keeps them coupled over a persisted project.
